\section{问题分析}
\label{sec:analysis}

\subsection{总体思路}

本题的核心是“把算力花在哪里”。四个问题分别对应一条完整决策链上的四个环节：\textbf{度量}（数据好不好、怎么配）$\rightarrow$ \textbf{建模}（数据与规模如何共同决定 Loss）$\rightarrow$ \textbf{决策}（预算约束下如何分配）$\rightarrow$ \textbf{预测}（Loss 如何转化为能力、能力前沿走向何处）。附件 A、B、C 分别来自不同实验体系，Loss 的量纲与模型族并不相同，直接拼接回归会把“数据来源差异”误当作“资源效应”。因此本文坚持\textbf{分源估计、接口传递}的原则：每一问只用本问对应的附件估计参数，跨问题只传递含义明确的中间量（质量基线 $Q_0$、参考配比 $\boldsymbol p_{\mathrm{ref}}$、配比响应 $R(\boldsymbol p)$、广义标度律参数、最优 Loss），并在接口处提出可检验的假设。总体技术路线如图~\ref{fig:overall}所示。

\begin{figure}[htbp]
\centering
\resizebox{\linewidth}{!}{% 全文技术路线图（TikZ，配色与插图统一）
\begin{tikzpicture}[
  font=\footnotesize,
  >=Stealth,
  box/.style={draw=#1, fill=#1!8, rounded corners=2pt, line width=0.6pt,
              align=center, minimum height=12mm, text width=31mm, inner sep=2pt, font=\scriptsize},
  qlab/.style={draw=#1, fill=#1, text=white, rounded corners=2pt, align=center,
               minimum height=12mm, text width=6mm, font=\scriptsize\bfseries},
  head/.style={font=\small\bfseries, text=pblue},
  arr/.style={->, line width=0.6pt, draw=pgray},
  link/.style={->, line width=0.8pt, draw=pred},
]
\def\dx{40mm}\def\dy{19mm}
% 列标题
\node[head] at (1*\dx-4mm, 0.9*\dy) {数据};
\node[head] at (2*\dx-4mm, 0.9*\dy) {模型与方法};
\node[head] at (3*\dx-4mm, 0.9*\dy) {关键输出};
\node[head] at (4*\dx-4mm, 0.9*\dy) {检验与分析};

% ---------------- 问题一
\node[qlab=pblue] (q1) at (0.25*\dx, 0) {问\\题\\一};
\node[box=pblue] (d1) at (1*\dx-4mm, 0) {A1--A3 质量信号\\A4--A15 配比实验\\A16/A18 映射与文本};
\node[box=pblue] (m1) at (2*\dx-4mm, 0) {分块稳定 CRITIC 赋权\\Huber 冲突消解\\Scheff\'e--Ridge 混料模型};
\node[box=pblue] (o1) at (3*\dx-4mm, 0) {领域质量 $q_d$、$Q_0$\\近优配比 $\boldsymbol p^\star$\\配比响应 $R(\boldsymbol p)$};
\node[box=pblue] (v1) at (4*\dx-4mm, 0) {抽样集/扩展集对照\\1M/60M/1B 检验集\\10B/70B 外推};

% ---------------- 问题二
\node[qlab=pgreen] (q2) at (0.25*\dx, -\dy) {问\\题\\二};
\node[box=pgreen] (d2) at (1*\dx-4mm, -\dy) {B1 规模轨迹\\B7 质量网格\\A6--A11 三规模配比};
\node[box=pgreen] (m2) at (2*\dx-4mm, -\dy) {留一规模选模\\成块交叉验证\\配比尺度不变性检验};
\node[box=pgreen] (o2) at (3*\dx-4mm, -\dy) {广义标度律\\弹性与边际效用\\质量—规模替代条件};
\node[box=pgreen] (v2) at (4*\dx-4mm, -\dy) {数据坍缩检验\\B2/B4/B5/B10 验证\\B8 压力测试};

% ---------------- 问题三
\node[qlab=porange] (q3) at (0.25*\dx, -2*\dy) {问\\题\\三};
\node[box=porange] (d3) at (1*\dx-4mm, -2*\dy) {三档预算 $C$\\三类质量成本 $g$\\C7 五档上下文};
\node[box=porange] (m3) at (2*\dx-4mm, -2*\dy) {配比可分性定理\\解析降维 + 多起点\\KKT 活跃集};
\node[box=porange] (o3) at (3*\dx-4mm, -2*\dy) {最优 $(N^*,D^*,Q^*)$\\临界上下文 30000\\三阶段结构性转移};
\node[box=porange] (v3) at (4*\dx-4mm, -2*\dy) {二维/三维复核\\参数 bootstrap\\234 个敏感性情景};

% ---------------- 问题四
\node[qlab=pred] (q4) at (0.25*\dx, -3*\dy) {问\\题\\四};
\node[box=pred] (d4) at (1*\dx-4mm, -3*\dy) {C1/C2 能力面板\\C3/C4 历史与算力\\C6 桥接、C8 逐任务};
\node[box=pred] (m4) at (2*\dx-4mm, -3*\dy) {Loss--能力单调桥接\\0.90 分位动力学前沿\\Shapley 贡献分解};
\node[box=pred] (o4) at (3*\dx-4mm, -3*\dy) {规模/非规模贡献\\三种算力情景\\12/24 月能力前沿};
\node[box=pred] (v4) at (4*\dx-4mm, -3*\dy) {来源族留一验证\\滚动起点回测\\联合块 bootstrap};

% 行内箭头
\foreach \i in {1,2,3,4}{
  \draw[arr] (d\i) -- (m\i);
  \draw[arr] (m\i) -- (o\i);
  \draw[arr,dashed] (o\i) -- (v\i);
}
% 跨问题接口
\draw[link] (o1.south) -- node[right,font=\scriptsize,text=pred,xshift=1mm]{$Q_0,\ \boldsymbol p^\star,\ R(\boldsymbol p)$} (o2.north);
\draw[link] (o2.south) -- node[right,font=\scriptsize,text=pred,xshift=1mm]{广义标度律} (o3.north);
\draw[link] (o3.south) -- node[right,font=\scriptsize,text=pred,xshift=1mm]{最优 Loss $L^*$} (o4.north);
\end{tikzpicture}
}
\caption{全文技术路线：四个问题的数据来源、模型方法、关键输出与检验关系}
\label{fig:overall}
\end{figure}

\subsection{问题一的分析}

问题一包含两条统计单位不同的任务。\textbf{质量评价}以文本为单位：22 项指标兼有连续统计量、分类概率与列表型评分，量纲不一、方向各异，且长度、句数等结构指标并非越大越好，需要先做“同向化 + 适中型效用”处理；指标间信息高度重叠（如三类相似度指标），直接等权会重复计权，因此采用\textbf{按语义分块、块内客观赋权、块间等权}的层次综合评价。当“内容好而形式差”一类矛盾出现在同一文本上时，简单平均会被单一极端维度拉偏，需要先用块间极差识别冲突，再用对离群块不敏感的 Huber 稳健估计消解。

\textbf{配比建模}以训练配方为单位：17 个领域份额满足单纯形约束 $\sum_ip_i=1$，普通回归会出现完全共线。混料试验设计中的 Scheff\'e 多项式\cite{scheffe1958mixtures,cornell2002mixtures}天然满足该约束，其二阶交互项还能直接刻画“两个领域组合是否产生额外收益”。153 个特征相对 512 组训练配方偏多，需要 Ridge 正则化；模型在 1M、60M、1B 三个实测规模上检验，再用 10B、70B 估算表讨论外推。

\subsection{问题二的分析}

经典标度律只含 $N,D$。本问要把质量 $Q$ 与配比 $\boldsymbol p$ 纳入，但附件中不存在同时改变四个因素的联合实验：B1 只变 $(N,D)$，B7 在固定 $(N,D)$ 网格上变 $Q$，附件 A 只变 $\boldsymbol p$。因此广义标度律必须采用\textbf{可分结构}，使每一项都能由一类数据单独识别，再通过可检验的假设拼合：(i) 质量项在 $Q=1$ 时消失，模型退化为经典形式；(ii) 配比项以\textbf{相对可约损失}的形式作用，其强度不随规模变化——这一点可以在附件 A 的 1M、60M、1B 三个规模上直接检验。在此基础上，边际效用与弹性、质量—规模等价增幅都可以写成闭式，从而推导出“质量提升能否由规模扩张替代”的可计算条件。

\subsection{问题三的分析}

问题三是一个带非线性预算约束的连续优化问题。由于 Loss 关于 $D$ 严格递减，最优解必然用尽预算，可由预算方程解析消去 $D$，把三维问题降为 $(\log N,Q)$ 上的二维有界问题；若广义标度律中配比项为乘性结构，则最优 $(N^*,D^*,Q^*)$ 与配比无关，配比可以与规模—质量决策\textbf{分离}求解。“结构性转移”需要严格定义：本文以 KKT 条件中活跃约束集合的改变定义主要转移（如质量投入从“不启动”到“启动”再到“饱和”），以配置弹性的断点定义次要转移，并用参数 bootstrap 检验转移的稳定性。上下文长度只改变单位 Token 的有效训练成本，注意力与训练开销之比 $\eta L_{\mathrm{ctx}}/6$ 与预算无关，临界长度可解析给出。

\subsection{问题四的分析}

问题四从 Loss 尺度转向下游 Benchmark 能力。首先需要统一能力口径（六任务官方归一化均分）、开源口径（开放权重且可追溯）、模型类型（区分 pretrained 与 chat/finetuned）与时间轴（提交日期）。其次，“规模扩张”与“非规模技术进步”在历史数据中同时发生，需要在同一模型中同时控制训练算力与时间状态：本文以 $0.90$ 条件分位刻画能力前沿，以局部线性趋势描述非规模状态的演进，再用两因素 Shapley 值把首末时点的前沿提升无序依赖地分解为两部分。最后，以 C6 的 Loss--Benchmark 配对建立单调桥接，把问题三的最优 Loss 映射为能力，并在算力增速保留 $100\%$、$50\%$、$25\%$ 三种情景下外推 12、24 个月的能力前沿，用分层 bootstrap 同时传播桥接、前沿与算力趋势三类不确定性。
